\section{Prime Factorization Changes}
\label{sec:factorization}

\subsection{Factorization Table}

\begin{table}[h]
\centering
\caption{Prime factorisation of 2020 and 2030 seat counts for states
         with projected changes. ``Depth'' is the depth of the \ar{}
         bisection tree.}
\label{tab:factorization}
\begin{tabular}{lcccccc}
\toprule
State & 2020 $k$ & 2020 split & Depth & 2030 $k$ & 2030 split & Depth \\
\midrule
Texas & 38 & $19\times2$ & 2 & 41 & $20+21$ & 4 \\
California & 52 & $13\times4$ & 3 & 50 & $5\times10$ & 3 \\
Florida & 28 & $7\times4$ & 3 & 29 & $14+15$ & 3 \\
New York & 26 & $13\times2$ & 2 & 25 & $5\times5$ & 2 \\
Arizona & 9 & $3\times3$ & 2 & 10 & $5\times2$ & 2 \\
Georgia & 14 & $7\times2$ & 2 & 15 & $5\times3$ & 2 \\
N. Carolina & 14 & $7\times2$ & 2 & 15 & $5\times3$ & 2 \\
Illinois & 17 & $8+9$ & 4 & 16 & $2\times8$ & 4 \\
Michigan & 13 & $6+7$ & 4 & 12 & $3\times4$ & 3 \\
Pennsylvania & 17 & $8+9$ & 4 & 16 & $2\times8$ & 4 \\
\bottomrule
\end{tabular}
\end{table}

Depth counts recursive split levels under the implemented
\texttt{split\_prescription}.  A large prime therefore has depth greater than
one because it first falls back to a binary split and then recurses on the two
composite or smaller-prime sides.

\subsection{Texas: The Most Dramatic Change}

Texas goes from $38 = 2 \times 19$ to $41$ (prime).
Under the current implementation, Texas first splits into 19 two-seat
containers.  Each two-seat container then performs a direct two-way split.
The large factor is used first so that states sharing a large factor can reuse
the same top-level partition.

Under the 2030 scenario tree, 41 is prime and greater than 3, so the root
uses the large-prime fallback: a binary split into 20-seat and 21-seat
subregions.  The 20 side then begins with a $5\times4$ split, while the
21 side begins with a $7\times3$ split.  The change from a 19-container
top level to a 20/21 binary top level means that the geographic organization
of the state's districts is fundamentally different.

\paragraph{Expected boundary movement.}
With $38 \to 41$ districts, the target district population decreases
from $P_{\rm state}/38$ to $P_{\rm state}/41$ (approximately 7.5\%
smaller per district).
Additionally, the tree structure changes.
We expect boundary movement of approximately 15--25\% of tracts
(compared to 3--5\% for seed changes within the current tree structure).

\subsection{California: Structural But Not Dramatic}

California goes from $52 = 2^2 \times 13$ to $50 = 2 \times 5^2$.
Both are composite, but the factorisation paths differ:
\begin{itemize}
\item $52$: split into 13 four-seat containers; each 4 splits into two
      two-seat containers, then each two-seat container splits directly.
      Tree depth: 3.
\item $50$: split into 5 ten-seat containers; each 10 splits into five
      two-seat containers, then each two-seat container splits directly.
      Tree depth: 3.
\end{itemize}

Both have depth 3, but the intermediate splitting differs.
At the top level, the $52$ tree makes 13 containers of 4 while the $50$ tree
makes 5 containers of 10.
This is a structural change but less dramatic than the Texas prime transition.

\subsection{Florida: The Most Structurally Significant Change}

Florida goes from $28 = 2^2 \times 7$ (composite: seven four-seat containers)
to $29$ (prime: binary $14:15$ fallback).

The seat gain changes the tree type from composite to prime-fallback,
the reverse of the prime-to-composite transitions in Illinois, Michigan,
and Pennsylvania.
This is Florida's most structurally significant scenario: 29 is prime,
so the 29-seat tree begins with a $14:15$ split at the root.
By contrast, a hypothetical $30 = 2 \times 3 \times 5$ would allow a
hierarchical tree ($15 \to 5 \to 3$ or $10 \to 6 \to 2$ depending on
factor order), which can better preserve county boundaries because the
intermediate bisections are aligned with the state's geographic structure.

The composite-to-prime transition ($28 \to 29$) is structurally analogous
to the Texas case ($38 \to 41$) but on a smaller scale.
Florida's case is the most structurally significant change among the
seat-gaining states because the alternative 30-seat composite tree would
preserve county boundaries substantially better than the 29-seat prime
fallback.

\subsection{Illinois, Michigan, Pennsylvania: Prime to Composite}

These states currently have prime seat counts ($17, 13, 17$) and will
move to composite ($16, 12, 16$).
A large prime $k$ has a binary fallback tree.
A composite $k$ has a largest-prime-first hierarchical tree.

This is the reverse of the Texas transition: going from a prime fallback to a
composite largest-prime-first tree.  Boundary movement from restructuring the
tree can still be substantial, even when both prescriptions are deterministic.

\begin{table}[h]
\centering
\caption{Qualitative characterisation of the 2030 tree-structure changes.}
\label{tab:tree-changes}
\begin{tabular}{lcc}
\toprule
State & 2020 tree type & 2030 tree type \\
\midrule
Texas & 19 containers of 2 & Binary 20/21 fallback \\
California & 13 containers of 4 & 5 containers of 10 \\
Florida & 7 containers of 4 & Binary 14/15 fallback \\
New York & Two-level (2-13) & Two-level (5-5) \\
Arizona & Two-level (3-3) & Two-level (2-5) \\
Georgia/NC & Two-level (2-7) & Two-level (3-5) \\
Illinois/PA & Binary 8/9 fallback & Four-level ($2^4$) \\
Michigan & Binary 6/7 fallback & Three-level ($3 \times 4$) \\
\bottomrule
\end{tabular}
\end{table}
